\documentclass[11pt]{article}
\usepackage[utf8]{inputenc}
\usepackage{amsmath}
\usepackage{amsfonts}
\usepackage{amssymb}
\usepackage{bm}
\usepackage{geometry}
\usepackage{xcolor}
\usepackage{booktabs}

\geometry{margin=0.8in}

\title{EM-FDR-Bagging: Robust Estimation of the Connectivity Matrix\\
from Multi-State Spiking Data\\
\large A bagged, scramble-calibrated wrapper around PRISM-EM}
\author{Jan Balewski, NERSC}
\date{\today}

\begin{document}

\maketitle

\begin{abstract}
This document describes \textbf{EM-FDR-Bagging}, a statistical wrapper around
the PRISM-EM fitter that aims to recover the shared connectivity matrix
\(A\in\mathbb{R}^{N\times N}\) of a non-stationary, multi-state Poisson GLM
\emph{robustly}, i.e.\ with explicit control of false edges and with
per-edge uncertainty estimates, on real data where no ground truth exists.
The method combines three ideas. First, \emph{block bagging}: the recording
is resampled into multiple overlapping bags built from spaced contiguous
time blocks, and PRISM-EM is fit independently on each bag. Second, a
\emph{circular-shift null}: within each bag the fitted state sequence is
frozen and the connectivity and state biases are re-estimated many times on per-neuron
time-scrambled spikes, producing a null distribution of edge weights that
have been passed through exactly the same estimation machinery as the real
weights. Third, \emph{per-row pooled significance with cross-bag stability
selection}: edges are tested against a row-pooled null and retained only if
they are selected reliably across bags. The procedure yields, for every
surviving edge, a selection frequency, a bagged point estimate, a bootstrap
SD, and a null-referenced effect size. We treat PRISM-EM as a
black box exposed through two entry points and specify everything needed for
an implementer who has the PRISM-EM code in hand.
\end{abstract}

\tableofcontents
\newpage

% ===============================================================
\section{Scope and Inputs}
% ===============================================================

The input is a single spike-train recording of \(N\) neurons over \(T\)
discrete time bins of width \(\Delta t\), in the format consumed by
PRISM-EM (\texttt{spikesData/<dataName>.spikes.npz}). The target regime,
which sizes every numerical choice below, is:
\(N=50\)–\(500\) neurons; recordings of \(10\)–\(60\) minutes at
\(\Delta t=10\,\text{ms}\), so \(T\) ranges from \(6\times10^{4}\) to
\(3.6\times10^{5}\) bins; per-neuron firing rates of \(1\)–\(5\,\text{Hz}\);
state-switching rates of \(5\)–\(15\) transitions per minute; and a small
number of latent states, typically \(M\approx3\). Crucially, the method is
designed for application to \emph{real} data: it never references ground
truth, and every quantity it reports is computed from the data alone.

The deliverable is not a single point estimate of \(A\) but a vetted edge
list. For each off-diagonal entry \((i,j)\) that survives the procedure we
report a point estimate, a statistical uncertainty, a stability (selection)
frequency, and an effect size relative to the null. Self-couplings
\(A_{ii}\) are excluded from testing throughout.

% ===============================================================
\section{The PRISM-EM Workhorse and Its Two Entry Points}
\label{sec:workhorse}
% ===============================================================

EM-FDR-Bagging treats PRISM-EM as a callable. We summarize only what the
wrapper relies on; the full model, objective, and defaults are documented
separately in the PRISM-EM specification.

\subsection{What PRISM-EM fits}

PRISM-EM fits a non-stationary Poisson GLM in which all \(M\) latent states
share one connectivity matrix \(A\) but carry state-specific bias vectors
\(B_m\). With \(c_t\) the simplex coefficients at bin \(t\) and
\(B_t=\sum_m c_{mt}B_m\), the rate is
\(\mu_t=\exp\!\big(\min(A Y_{t-1}+B_t,\,\eta_{\text{clip}})\big)\,\Delta t\)
and \(Y_t\sim\text{Poisson}(\mu_t)\). Training is block coordinate descent:
an E-step infers \(c_{1:T}\) by simplex-projected gradient descent and
Viterbi-decodes them to a hard one-hot state sequence
\(\hat S_{1:T}\); an M-step then updates \(A,\mathbf{B}\) by Adam under the
frozen one-hot assignment, with off-diagonal \(L_1\) shrinkage (strength
\(\lambda_3\)), a spectral-radius projection
\(\rho(A)\le\rho_{\max}\), and a learning-rate decay. The
shrinkage, projection, and decay each switch on only after a configurable
number of EM iterations. The fitted \(A\) is therefore the output of a
specific, deterministic-in-configuration pipeline of regression plus
shrinkage plus spectral projection --- a fact the null construction below
depends on.

\subsection{Entry point 1: full fit}
\label{sec:ep1}

\begin{quote}
\(\textsf{full\_fit}(\text{spikes},\,M,\,\theta_{\text{EM}})
\;\longrightarrow\;(\hat A,\ \hat S_{1:T},\ \hat{\mathbf B},\ \text{history})\)
\end{quote}

This is the ordinary PRISM-EM run: full E/M alternation on the supplied
spikes, returning the fitted connectivity \(\hat A\), the decoded one-hot
state sequence \(\hat S_{1:T}\), the per-state biases \(\hat{\mathbf B}\),
and the training history. Here \(\theta_{\text{EM}}\) collects all PRISM-EM
hyperparameters (number of EM iterations, M-step epochs, learning rates,
\(\lambda_2,\lambda_3,\rho_{\max}\), the delay schedules, etc.).

\subsection{Entry point 2: M-step-only refit with locked state sequence}
\label{sec:ep2}

\begin{quote}
\(\textsf{mstep\_locked}(\,Y_{\text{prev}},\,Y_{\text{curr}},\,\hat S^{\text{lock}},\,
A^{\text{init}},\,\hat{\mathbf B}^{\text{init}},\,\theta_{\text{M}}\,)
\;\longrightarrow\;(A,\mathbf B)\)
\end{quote}

This entry point runs \emph{only} the M-step. It accepts spike pairs
\((Y_{\text{prev}},Y_{\text{curr}})\), an externally supplied and
\emph{frozen} state sequence \(\hat S^{\text{lock}}\) (used directly as the
one-hot assignment, with no E-step and no Viterbi re-decoding), initial
parameters \((A^{\text{init}},\hat{\mathbf B}^{\text{init}})\), and returns
the re-estimated connectivity \(A\) and the refit state biases
\(\mathbf B\). In the null workflow the initial \(A\) and \(\mathbf B\) are
not copied from the real time-ordered fit. Instead, for every scramble they
are initialized freshly from the scrambled data, or randomly, according to
\texttt{null\_A\_init} and \texttt{null\_B\_init}. Re-fitting
\(\mathbf B\) in the scrambled context lets the null absorb state-specific
baseline effects, so the test of \(A\) is not made artificially harsh by a
bias mismatch; fresh \(A\) initialization avoids seeding the null from the
real edge pattern.

\paragraph{Locked-state time indexing.} For each lag-1 training pair
\((Y_{t-1},Y_t)\), the locked state used by the M-step is
\(\hat S^{\text{lock}}_{t-1}\), i.e.\ the state attached to the input/history
bin. Equivalently, if the spike pairs have length \(T-1\), the state
covariate rows are formed from \(\hat S^{\text{lock}}_{0:T-2}\), converted
to one-hot vectors. This convention is part of the model definition for
lag-1 fitting and must be used identically for the real time-ordered fit and
for every scrambled locked refit.

The locked M-step is configured to reproduce the real
pipeline exactly: the same \(L_1\) strength \(\lambda_3\), the same spectral
projection \(\rho_{\max}\), and the same learning-rate decay as the
corresponding \textsf{full\_fit}. Because there is no E-step here, the EM
iteration counters that gate PRISM-EM's delayed schedules are not
meaningful; instead all schedule-gated operations (shrinkage, spectral
projection, decay) are \textbf{active from the first M-step epoch}. The
rationale is that the null estimate must traverse the same shrinkage and
constraint geometry as the converged real estimate, which lives in the
late-EM regime where all schedules have already fired.

\paragraph{Implementation note (prerequisite).} The current PRISM-EM code
inlines the M-step inside the EM loop and feeds it the one-hot sequence
produced internally by the E-step. To support entry point~2, the M-step
must be exposed as a standalone callable that consumes an externally
supplied \(\hat S^{\text{lock}}\) and externally supplied spike pairs. This
includes making the lag-1 state alignment explicit as
\(\hat S^{\text{lock}}_{t-1}\) for pair \((Y_{t-1},Y_t)\), and allowing
the caller to provide explicit initial \(A\) and \(\mathbf B\). This is a
small refactor of existing code rather than new modeling work; the rest of
this document assumes both entry points exist with the signatures above.

% ===============================================================
\section{Method Overview}
% ===============================================================

EM-FDR-Bagging has two stages, and the document is organized to match the
intended software structure.

\textbf{Stage~(a), per-bag fitting} (Sections~\ref{sec:bagging}--\ref{sec:bagout}).
For a single bag: assemble a resampled time series from spaced blocks, run
\textsf{full\_fit} to obtain the real \((\hat A^{(b)},\hat S^{(b)})\), then run
\(P\) scramble refits through \textsf{mstep\_locked} with the state sequence
frozen, obtaining a set of null connectivity matrices. Write everything to
disk. This stage is one self-contained program; it is run \(N_{\text{bag}}\)
times, once per bag, as independent jobs.

\textbf{Stage~(b), aggregation and pruning} (Section~\ref{sec:aggregate}).
Read all per-bag outputs from a directory, build per-row pooled null
distributions, perform per-bag selection and cross-bag stability selection,
and emit the final vetted edge list with uncertainties.

The symbols introduced for the wrapper, kept distinct from PRISM-EM's
\(N,M,T\), are collected in Table~\ref{tab:symbols}.

\begin{table}[h]
\centering
\begin{tabular}{lll}
\toprule
Symbol & Meaning & Nominal value \\
\midrule
\(N_{\text{blk}}\) & blocks per bag & \(10\) \\
\(L_{\text{blk}}\) & block length (s) & \(60\) \\
\(L_{\text{blk}}^{\text{bin}}\) & block length (bins) \(=L_{\text{blk}}/\Delta t\) & \(6000\) \\
\(R_{\text{time}}\) & time range for block starts (s) & e.g.\ \([1000,1600]\) \\
\(N_{\text{bag}}\) & number of bags & \(5\)--\(10\) \\
\(P\) & scrambles per bag & \(4\)--\(8\) \\
\(\delta_{\text{blk}}^{\min}\) & min separation between block starts (s) & \(1\) \\
\(\delta_{\text{roll}}^{\min}\) & roll exclusion-zone half-width (s) & \(2\) \\
\(q\) & per-bag per-row selection quantile level & e.g.\ \(0.99\) \\
\(\pi_{\text{thr}}\) & cross-bag selection-frequency threshold & \(0.6\)--\(0.8\) \\
\(M,N,T\) & states, neurons, bins (PRISM-EM) & \(\approx3\); \(50\)--\(500\); \(6\!\times\!10^4\)--\(3.6\!\times\!10^5\) \\
\bottomrule
\end{tabular}
\caption{Symbols introduced by EM-FDR-Bagging and their nominal values.
Every numeric choice in the text is written as ``symbol = value'' so it can
be adjusted.}
\label{tab:symbols}
\end{table}

% ===============================================================
\section{Stage (a.1): Block Bagging}
\label{sec:bagging}
% ===============================================================

A bag is a surrogate recording assembled by concatenating
\(N_{\text{blk}}=10\) contiguous blocks, each of length
\(L_{\text{blk}}=60\,\text{s}\) (i.e.\ \(L_{\text{blk}}^{\text{bin}}=6000\)
bins at \(\Delta t=10\,\text{ms}\)), drawn from a designated time range
\(R_{\text{time}}\) of the original recording. The reassembled bag has length
\(N_{\text{blk}}\times L_{\text{blk}}^{\text{bin}}\) bins (here
\(6\times10^4\) bins, i.e.\ \(10\) minutes), which is the sequence PRISM-EM
fits.

\paragraph{Block placement.} Block start positions are bin indices drawn
from the time range. The range \(R_{\text{time}}\) is given in seconds and
must leave room for a full block, so admissible starts (in bins) lie in
\([\,R_{\text{time}}^{\text{start}}/\Delta t,\;
R_{\text{time}}^{\text{end}}/\Delta t - L_{\text{blk}}^{\text{bin}}\,]\).
Within that range a start may fall on essentially any bin, so the number of
candidate starts is large (tens of thousands), and exact collisions between
two chosen starts are vanishingly unlikely. The genuine hazard is not
identical starts but \emph{near-identical, heavily overlapping} blocks: two
starts one bin apart share all but one bin of content and would inject a
near-duplicate stretch into the bag. We therefore impose a minimum
separation \(\delta_{\text{blk}}^{\min}=1\,\text{s}\) (\(100\) bins) between
any two chosen starts and sample by rejection: draw \(N_{\text{blk}}\)
starts uniformly from the admissible range, reject and redraw any draw that
violates the minimum-separation constraint, and try up to
\texttt{max\_block\_draw\_trials}=30 times for each new block. If all 30
attempts fail, accept the 31st draw regardless of the separation constraint
and mark the violation in the job metadata. This bounded retry policy avoids
silent infinite loops while still giving the ordinary case many chances to
avoid near-duplicates.

\paragraph{Reference chronological bag.} The special job
\texttt{bag\_idx=0} is not randomly resampled. Its block starts follow the
input-data timeline: block 0 begins at the start of \texttt{--time\_range\_sec},
block 1 begins one block length later, and so on. This bag is used as the Stage~(b)
reference for time-series fields, so state-assignment plots from the
aggregate output remain aligned to the original simulation timeline when
ground truth is available. All \texttt{bag\_idx>0} jobs use the random
block-bootstrap sampler described above.

\paragraph{Blocks are non-circular.} Each block is a literal contiguous slice
of the recording, clipped at the window end; no wrap-around is used in block
extraction. Lag structure within a block is preserved exactly, so the first
lag pair \((Y_0,Y_1)\) at a block boundary is the only mild discontinuity,
and with \(L_{\text{blk}}^{\text{bin}}=6000\) bins per block the boundary
fraction is negligible. By default these cross-block boundary pairs are
included rather than masked, because the assembled bag is passed to
\textsf{full\_fit} as an ordinary spike sequence and the number of such
pairs is only \(N_{\text{blk}}-1\) out of \(T_b-1\). The realized block
start offsets and the assembled-sequence boundary pair indices are stored
in the per-bag job record so their influence can be audited, or a stricter
masked-pair variant can be evaluated later.

\paragraph{Why overlapping bags.} Bags are deliberately allowed to overlap
in source time across the \(N_{\text{bag}}\) runs; this is a block bootstrap
on a single finite recording. Overlap induces positive correlation between
bags, which has consequences for the uncertainty estimates
(Section~\ref{sec:errors}) and the stability bound
(Section~\ref{sec:aggregate}); these are stated honestly where they bite.
The overlap can be reduced --- and its impact studied --- by lengthening the
total recording while holding \(N_{\text{blk}}\) small, which the parameter
layout supports directly.

% ===============================================================
\section{Stage (a.2): Real Fit and the Circular-Shift Null}
\label{sec:null}
% ===============================================================

\subsection{Real fit on the bag}

For bag \(b\), call entry point~1 on the assembled bag:
\[
(\hat A^{(b)},\ \hat S^{(b)}_{1:T_b},\ \hat{\mathbf B}^{(b)})
=\textsf{full\_fit}(\text{bag}_b,\,M,\,\theta_{\text{EM}}),
\qquad T_b=N_{\text{blk}}L_{\text{blk}}^{\text{bin}}.
\]
This \(\hat A^{(b)}\) is the bag's real connectivity estimate; the frozen
state sequence \(\hat S^{(b)}\) is reused unchanged for all scrambles below.

\subsection{The null: per-neuron independent circular shifts}

The null hypothesis for an edge is ``no directed lag-1 coupling from neuron
\(j\) to neuron \(i\) beyond what the bias/state structure already explains.''
We realize it by destroying cross-neuron temporal alignment while preserving
each neuron's own marginal statistics and the state timeline. Concretely,
for scramble \(p=1,\dots,P\) we draw an independent integer circular shift
\(\delta_n^{(p)}\) for each neuron \(n\) and roll that neuron's spike train:
\[
\tilde Y^{(p)}_{t,n} = Y_{(t-\delta_n^{(p)})\bmod T_b,\ n}.
\]
Because every neuron is shifted by a different random amount, the precise
lag-1 relationship that \(A\) captures is broken, while each neuron's firing
rate, burstiness, and count distribution are untouched (a circular roll is a
pure permutation of a neuron's own time axis). The state sequence
\(\hat S^{(b)}\) is \emph{not} rolled: it is held fixed so that the null asks
specifically about coupling, not about state structure.

\paragraph{Two-sided exclusion zone on each shift.} A shift near \(0\) leaves
the data almost unshifted and a shift near \(T_b\) wraps back to almost
unshifted; both leave the real lag structure largely intact and would
contaminate the null toward the alternative. We therefore restrict each
shift to
\[
\delta_n^{(p)}\in
\big[\,\delta_{\text{roll}}^{\min,\text{bin}},\ \
T_b-\delta_{\text{roll}}^{\min,\text{bin}}\,\big],
\qquad
\delta_{\text{roll}}^{\min,\text{bin}}=\delta_{\text{roll}}^{\min}/\Delta t,
\]
with \(\delta_{\text{roll}}^{\min}=2\,\text{s}\) (\(200\) bins). The
exclusion is two-sided by construction: both the small-shift and the
near-full-wrap small-effective-shift regions are forbidden.

\paragraph{No pairwise separation among neurons.} The shifts of different
neurons need not be separated from one another. The null only requires that
each neuron's shift is individually non-trivial (the two-sided zone above)
and that shifts are drawn independently across neurons; a coincidence in
which two neurons happen to share the same shift is harmless to null
validity. This matters in practice: with a \(10\)-minute bag and a
\(2\,\text{s}\) separation there would be only \(\sim300\) ``slots,'' yet a
bag may contain up to \(500\) neurons, so any attempt to enforce pairwise
distinct shifts would be infeasible by pigeonhole. We deliberately avoid
that requirement. A lightweight \texttt{max\_retries} guard on the per-neuron
draw is retained only to handle a pathological window, not to enforce mutual
spacing.

\subsection{Null re-estimation through the locked M-step}

For each scramble \(p\) we recompute connectivity with the state sequence
frozen, via entry point~2:
\[
(A^{\text{null},(b,p)},\mathbf B^{\text{null},(b,p)})
=\textsf{mstep\_locked}\!\big(\tilde Y^{(p)}_{\text{prev}},\ \tilde Y^{(p)}_{\text{curr}},\
\hat S^{(b)},\ \hat{\mathbf B}^{(b)},\ \theta_{\text{M}}\big).
\]
The locked state attached to pair
\((\tilde Y^{(p)}_{t-1},\tilde Y^{(p)}_t)\) is
\(\hat S^{(b)}_{t-1}\). The state sequence itself remains the decoded
sequence from the real, time-ordered bag fit; only the spike counts entering
the null M-step are circularly shifted. The null initialization is fresh for
each scramble: \(A^{\text{init},(b,p)}\) is either computed from the
scrambled spike train (\texttt{null\_A\_init=scrambled\_data}, the default) or drawn
randomly (\texttt{null\_A\_init=rand}), and
\(\mathbf B^{\text{init},(b,p)}\) is likewise initialized from scrambled
data (\texttt{null\_B\_init=scrambled\_data}, the default) or randomly
(\texttt{null\_B\_init=rand}). The real time-ordered
\(\hat A^{(b)}\) is deliberately not used to seed the null refit.
The configuration \(\theta_{\text{M}}\) matches the real pipeline exactly
(Section~\ref{sec:ep2}): identical \(\lambda_3\), identical \(\rho_{\max}\),
identical learning-rate decay, all active from the first epoch. This is the
crux of the calibration: the real \(\hat A^{(b)}\) and the null
\(A^{\text{null},(b,p)}\) are both outputs of the \emph{same} regression,
shrinkage, and spectral projection, applied to the same state timeline,
differing only in whether cross-neuron alignment was destroyed and in the
resulting refit biases. A
significance statement of the form ``this edge is larger than the same
machinery produces on scrambled data'' is therefore meaningful; using a raw,
unshrunk M-step for the null would compare against a differently biased
quantity and invalidate the test.

\subsection{How many scrambles are needed}
\label{sec:howmanyscrambles}

A useful efficiency follows from testing significance per row rather than
per edge (Section~\ref{sec:aggregate}). For row \(i\) we pool the
\((N-1)\) off-diagonal null weights of that row across all scrambles. A
\emph{single} scramble therefore already contributes \(\sim(N-1)\) null
samples to row \(i\)'s distribution --- of order \(199\) for \(N=200\). With
a modest \(P=4\)--\(8\) the pooled per-row null holds
\((N-1)\times P\approx800\)--\(1600\) samples, ample to resolve a \(1\%\)
tail. By contrast a \emph{per-edge} null would need on the order of hundreds
of scrambles just to populate a single edge's distribution, which is
infeasible at this budget. The two regimes are not equivalent: the
\((N-1)\) samples within one scramble are not independent (they share a
single scrambled realization of the data), so increasing \(P\) still helps
by adding genuinely independent null realizations and by averaging out
scramble-to-scramble variability. The recommendation is to keep \(P\) small
(\(4\)--\(8\)) but greater than one, precisely because per-row pooling
supplies the raw sample count while extra scrambles supply independence.

% ===============================================================
\section{Stage (a.3): Per-Bag Outputs}
\label{sec:bagout}
% ===============================================================

Each per-bag job writes a self-contained file to a shared results directory,
with a name carrying the bag index so that \(N_{\text{bag}}\) parallel jobs
do not collide. The file contains, at minimum: the real connectivity
\(\hat A^{(b)}\in\mathbb{R}^{N\times N}\); the stack of null connectivities
\(\{A^{\text{null},(b,p)}\}_{p=1}^{P}\) of shape \((P,N,N)\); the frozen
state sequence \(\hat S^{(b)}\); the realized block start indices and the
assembled-sequence boundary pair indices; any block-start constraint
violations caused by exhausting \texttt{max\_block\_draw\_trials}; the
realized per-neuron shifts (for reproducibility and for the data-reuse
study); the stack of refit null biases
\(\{\mathbf B^{\text{null},(b,p)}\}_{p=1}^{P}\) of shape \((P,M,N)\); the fitted
off-diagonal sparsity of \(\hat A^{(b)}\) (number and
fraction of surviving edges) as a sanity cross-check; and the PRISM-EM
training-history summary (final objective, iteration count, whether each
delayed schedule fired, final state occupancy). The assembled bag spike
train itself is not saved, because the ordered block starts are sufficient
to reconstruct it from the original input spikes if needed. Stage~(b)
consumes only these files and never re-runs PRISM-EM.

\paragraph{Per-bag program output naming.} One invocation of
\texttt{prism\_EM\_FDR\_Bags\_train3c.py} produces one bag file containing
the real time-ordered fit and all \(P\) scramble refits for that bag. The
filename is deterministic and includes \texttt{dataName}, \texttt{bagsTag},
and \texttt{bag\_idx}, for example
\texttt{<dataName>\_<bagsTag>.bag003.prismFDRbag.npz}. The real fit fields are written
in the same schema as the standalone \texttt{prism\_EM\_train3c.py} output,
so the usual baseline fit evaluation can be run on the real part of a bag
without a separate training job.

\paragraph{Per-bag driver runtime controls.} The bag driver inherits the
ordinary PRISM-EM training controls (\texttt{num\_states},
\texttt{num\_em\_iters}, \texttt{m\_epochs}, learning rates,
\(\lambda_2,\lambda_3,\rho_{\max}\), schedule delays, batch size, and
initialization modes for the real fit). Its additional controls are:
\texttt{bag\_idx}, \texttt{bagsTag}, \texttt{time\_range\_sec},
\texttt{num\_blocks}, \texttt{block\_len\_sec},
\texttt{min\_block\_start\_sep\_sec},
\texttt{max\_block\_draw\_trials}, \texttt{num\_scrambles},
\texttt{min\_roll\_shift\_sec}, \texttt{null\_A\_init}, and
\texttt{null\_B\_init}. There is no separate \texttt{null\_m\_epochs}
parameter: the locked null M-step always inherits the M-step epoch count and
optimizer configuration from the real time-ordered fit, with pruning,
spectral projection, and learning-rate decay active from the first locked
epoch. There is also no option to save the assembled bag spike train; the
ordered block starts are the reproducible representation of the bag.

\paragraph{Compute layout.} Each per-bag job is a single SLURM job on a
single node using \(4\times\)A100, matching the existing PRISM-EM launch
(\texttt{torchrun --standalone --nnodes=1 --nproc\_per\_node=4}). The
\(P\) scramble refits reuse the same node and the frozen \(\hat S^{(b)}\),
so they are cheap relative to the full fit. The \(N_{\text{bag}}\) jobs are
mutually independent and submitted in parallel.

% ===============================================================
\section{Stage (b): Aggregation, Significance, and Pruning}
\label{sec:aggregate}
% ===============================================================

Stage~(b) is a single postprocessing program pointed at the results
directory. It performs per-row null pooling, per-bag selection, and
cross-bag stability selection, then reports the vetted edge list with
uncertainties.

\subsection{Per-row pooled null and per-bag selection}
\label{sec:perrow}

We test significance \emph{per row}, exploiting the assumption that the
incoming edges to a given target neuron \(i\) are exchangeable under the null:
they share that target neuron's firing statistics, state occupancy, and
observation scale. This is a much milder assumption than per-edge normality.
The row convention is \(A_{ij}\) = effect of source neuron \(j\) on target
neuron \(i\); therefore this row-pooled test calibrates the incoming weights
to one target neuron. The separate biological interpretation of a source
neuron as broadly excitatory or inhibitory is read down a column of retained
edges, across the downstream targets it controls. The per-row exchangeability
assumption is what makes the small scramble budget viable
(Section~\ref{sec:howmanyscrambles}).

For bag \(b\) and row \(i\), pool the off-diagonal null weights across all
scrambles of that bag,
\[
\mathcal{N}^{(b)}_i
=\big\{\,A^{\text{null},(b,p)}_{ij}\ :\ j\neq i,\ p=1,\dots,P\,\big\},
\qquad |\mathcal{N}^{(b)}_i| = (N-1)\,P,
\]
and form a per-row threshold as the upper-\(q\) quantile of the magnitudes,
\[
\tau^{(b)}_i = \mathrm{Quantile}_{q}\!\big(\,\{|x| : x\in\mathcal{N}^{(b)}_i\}\,\big),
\qquad q = 0.99 \ \text{(nominal)}.
\]
An edge is \emph{selected in bag \(b\)} if its real magnitude exceeds the
row threshold,
\[
\text{sel}^{(b)}_{ij} = \mathbb{1}\!\big[\,|\hat A^{(b)}_{ij}| > \tau^{(b)}_i\,\big],
\qquad i\neq j.
\]
Equivalently, each real edge can be converted to a within-row \(p\)-value
against the pooled null and the Benjamini--Hochberg procedure applied within
the row; the quantile-threshold form above is the operational default. The
expected number of edges selected per bag and per row is recorded as a
diagnostic (it is the quantity \(q_\Lambda\) of the stability bound below);
it is \emph{measured}, never prescribed, so the algorithm discovers how many
edges it keeps rather than being told a target sparsity.

\subsection{Cross-bag stability selection}

A single bag's selection is noisy. We retain only edges that are selected
reliably across bags. Define the selection frequency
\[
\pi_{ij} = \frac{1}{N_{\text{bag}}}\sum_{b=1}^{N_{\text{bag}}}\text{sel}^{(b)}_{ij},
\]
and keep edge \((i,j)\) in the final network if
\(\pi_{ij}\ge\pi_{\text{thr}}\), with
\(\pi_{\text{thr}}=0.6\)--\(0.8\) (nominal). This is Meinshausen--B\"uhlmann
stability selection: the data are perturbed by resampling, the base selector
is run on each perturbation, and only consistently chosen variables survive.
For example, the implementation default \texttt{--sel\_prob 0.7} means that
an off-diagonal edge must be selected in at least \(70\%\) of the bags to
survive. The value assigned to a surviving edge is not weighted by the
selection frequency or by any bag quality score; it is the plain arithmetic
mean of \(\hat A^{(b)}_{ij}\) over the bags in which that edge passed the
per-bag null threshold. Bags where the edge did not pass the per-bag selector
do not enter this final \texttt{A\_hat} value, although the all-bag mean is
also saved as a diagnostic.

\paragraph{Two ways to set the knobs (one reported, one derived).} The
operational default exposes the pair \((q,\ \pi_{\text{thr}})\): a per-bag
per-row quantile level \(q\) and a cross-bag frequency threshold
\(\pi_{\text{thr}}\). A more compact, single-knob alternative collapses these
using the stability-selection error bound,
\[
\mathbb{E}[V]\;\le\;\frac{1}{2\pi_{\text{thr}}-1}\cdot\frac{q_\Lambda^{2}}{p_{\text{cand}}},
\]
where \(\mathbb{E}[V]\) is the expected number of false edges,
\(q_\Lambda\) is the \emph{measured} average number of edges selected per bag,
and \(p_{\text{cand}}=N(N-1)\) is the number of candidate off-diagonal edges.
In this mode the user fixes only a single target,
\texttt{target\_false\_edges}, interpreted as a target upper bound on
\(\mathbb{E}[V]\). The per-bag quantile is set loose and merely determines
\(q_\Lambda\), and \(\pi_{\text{thr}}\) is then \emph{derived} from the
bound rather than tuned. We document this route because it reduces the
user's choice to the expected number of false retained edges, but adopt the
explicit \((q,\pi_{\text{thr}})\) pair as the default for transparency. The
honest caveat applies to either route: the bound assumes independent
(exchangeable) bags, whereas our overlapping-block bags are positively
correlated, so a derived \(\pi_{\text{thr}}\) and any implied false-edge
control should be read as approximate and reported as such, not as a
guaranteed FDR.

% ===============================================================
\section{Per-Edge Uncertainty: Two Distinct Errors}
\label{sec:errors}
% ===============================================================

The design furnishes two genuinely different per-edge ``errors,'' which
answer different questions and must not be conflated.

\subsection{Null spread and effect size}

For row \(i\), the pooled null \(\mathcal{N}_i=\bigcup_b\mathcal{N}^{(b)}_i\)
has a mean and standard deviation,
\[
m^{\text{null}}_i=\operatorname{mean}\big(\mathcal{N}_i\big),
\qquad
s^{\text{null}}_i=\operatorname{SD}\big(\mathcal{N}_i\big),
\]
which characterize the \emph{noise floor}: how large an edge weight looks
purely by chance once cross-neuron alignment is destroyed and the same
machinery is applied. This is not the uncertainty of the estimate; it is the
spread of the null. It does, however, define a natural null-referenced
effect size,
\[
z_{ij}=\frac{\bar A_{ij}-m^{\text{null}}_i}{s^{\text{null}}_i},
\]
a per-edge signal-to-noise ratio against the scramble null, where
\(\bar A_{ij}\) is the bagged point estimate defined next.

\subsection{Estimation variability: the across-bag bootstrap SD}

The quantity one usually means by ``the statistical error of an edge'' is the
sampling variability of its fitted weight, and bagging estimates it directly.
Each bag yields a real fitted weight \(\hat A^{(b)}_{ij}\) on a different
resampled stretch of data; their spread across bags is a block-bootstrap
estimate of the estimator's sampling distribution. We report the bagged mean
and bootstrap SD,
\[
\bar A_{ij}=\frac{1}{B_{ij}}\sum_{b\in\mathcal{B}_{ij}}\hat A^{(b)}_{ij},
\qquad
\widehat{\mathrm{SD}}_{ij}
=\sqrt{\frac{1}{B_{ij}-1}\sum_{b\in\mathcal{B}_{ij}}\big(\hat A^{(b)}_{ij}-\bar A_{ij}\big)^{2}},
\]
where \(\mathcal{B}_{ij}\) is the set of bags entering the average and
\(B_{ij}=|\mathcal{B}_{ij}|\).

\paragraph{Which bags to average over.} A bag in which edge \((i,j)\) was
pruned to exactly zero by \(L_1\) contributes a structural zero, not a noisy
nonzero. Two summaries are therefore reported. The
\emph{selection-conditional} estimate averages only over bags where the edge
survived (\(\mathcal{B}_{ij}=\{b:\text{sel}^{(b)}_{ij}=1\}\)); it reflects the
typical magnitude when the edge is detected but is biased upward in
magnitude. The \emph{all-bag} estimate averages over every bag
(\(\mathcal{B}_{ij}=\{1,\dots,N_{\text{bag}}\}\), structural zeros included);
it is diluted toward zero in proportion to how often the edge is absent. We
recommend reporting the selection-conditional mean and bootstrap SD
\emph{together with} the selection frequency \(\pi_{ij}\), so that the dilution is visible
rather than hidden: \(\pi_{ij}\) and \(\bar A_{ij}\) together convey both how
often and how strongly an edge appears.

\paragraph{Honest caveat on calibration.} The bags are overlapping
block-bootstrap samples from one finite recording, not independent datasets.
Overlap makes bags more alike than independent draws, so the naive across-bag
SD \emph{underestimates} the true sampling error; with only
\(N_{\text{bag}}=5\)--\(10\) bags it is additionally coarse (few degrees of
freedom). Accordingly \(\widehat{\mathrm{SD}}_{ij}\) should be used as a
\emph{relative} reliability measure --- excellent for ranking edges and for
flagging unstable ones --- rather than as a calibrated frequentist standard
error supporting a strict confidence interval. The degree of underestimation
can be probed empirically by lengthening the recording and shrinking the
block overlap (Section~\ref{sec:bagging}), which is a deliberate affordance
of the design.

\subsection{Reported per-edge record}

For every edge that passes stability selection
(\(\pi_{ij}\ge\pi_{\text{thr}}\)) the pipeline emits
\[
\big(\,i,\ j,\ \pi_{ij},\ \bar A_{ij},\ \widehat{\mathrm{SD}}_{ij},\ z_{ij}\,\big),
\]
together with the all-bag mean for comparison and the row null summary
\((m^{\text{null}}_i,s^{\text{null}}_i)\). The selection frequency
\(\pi_{ij}\), the bagged point estimate, the bootstrap SD, and the
null-referenced effect size are complementary: an edge that is frequently
selected, large relative to its bootstrap SD, and far into the null tail is
trustworthy; an edge strong in one bag but rarely selected, or with bootstrap SD
comparable to its magnitude, is flagged as unreliable.

% ===============================================================
\section{Implementation and Runtime Arguments}
\label{sec:implementation}
% ===============================================================

This section records the concrete command-line interface used by the current
implementation. Stage~(a) produces one bag file per invocation. Stage~(b)
reads the ordered set of bag files and writes a single aggregate
\texttt{.prismEM.npz} file that can be opened by the existing
\texttt{prism\_EM\_eval3c.py} machinery.

\subsection{\texttt{prism\_EM\_FDR\_Bags\_train3c.py}: Stage (a)}

This program is run once per bag, normally as one multi-GPU job:
\[
\texttt{torchrun --standalone --nnodes=1 --nproc\_per\_node=4
./prism\_EM\_FDR\_Bags\_train3c.py ...}
\]
It assembles one block-resampled bag, performs the real time-ordered
PRISM-EM fit on that bag, runs the requested circular-shift locked-M-step
null refits, and writes one self-contained bag file.

\paragraph{Input and output naming.}
The input spike file is
\[
\texttt{<basePath>/spikesData/<dataName>.spikes.npz}.
\]
The output directory is \texttt{fdr\_out\_dir} if supplied, otherwise
\[
\texttt{<basePath>/prismFDR}.
\]
The output file is deterministic in \texttt{dataName}, \texttt{bagsTag}, and
\texttt{bag\_idx}:
\[
\texttt{<basePath>/prismFDR/<dataName>\_<bagsTag>.bag003.prismFDRbag.npz}
\]
for \(\texttt{bag\_idx}=3\) and default output directory. The bag index is
not only a label: it also enters the random seed stream, so rerunning the
same command with the same \texttt{bag\_idx} reproduces the same bag and
scramble shifts, while changing only \texttt{bag\_idx} produces a different
bag file. The one exception is \texttt{bag\_idx=0}: its block starts are
chronological by construction, while its scramble shifts still use the
bag-indexed random seed.

A single Stage~(a) bag can be evaluated directly by passing the bag stem to
\texttt{prism\_EM\_eval3c.py}; if \texttt{dataName} contains a
\texttt{bagNNN} token, the evaluator reads
\texttt{<basePath>/prismFDR/<dataName>.prismFDRbag.npz} instead of looking
under \texttt{prismFit}.

\paragraph{Required arguments.}
\begin{description}
  \item[\texttt{--basePath}] Root directory containing \texttt{spikesData/}
  and receiving \texttt{prismFDR/}. This is the same run directory used by
  the plain PRISM-EM trainer.
  \item[\texttt{--dataName}] Short dataset name. The program reads
  \texttt{<dataName>.spikes.npz} from \texttt{spikesData/}.
  \item[\texttt{--num\_states}] Number \(M\) of latent states in the
  switching Poisson GLM.
  \item[\texttt{--bag\_idx}] Integer bag index. It is used in the output
  filename and in the per-bag RNG seed. For \texttt{bag\_idx=0}, the bag
  blocks are chronological, starting at the time-range start and advancing
  by one block length per block; other bag indices draw random block starts.
  \item[\texttt{--bagsTag TAG}] Alphanumeric tag appended after
  \texttt{dataName} in Stage~(a) output names. If omitted or set to
  \texttt{None}, the program derives a deterministic four-character
  alphanumeric tag from the shared bagging/training configuration, excluding
  \texttt{bag\_idx}.
  \item[\texttt{--time\_range\_sec START END}] Time range, in
  seconds, from which block starts are drawn.
\end{description}

\paragraph{Bagging and scramble arguments.}
\begin{description}
  \item[\texttt{--fdr\_out\_dir}] Optional output directory for Stage~(a)
  bag files. Default: \texttt{<basePath>/prismFDR}.
  \item[\texttt{--num\_blocks}] Number of contiguous blocks concatenated
  into one bag. Default: \(10\).
  \item[\texttt{--block\_len\_sec}] Length of each contiguous block in
  seconds. Default: \(60\).
  \item[\texttt{--min\_block\_start\_sep\_sec}] Desired minimum separation
  between block starts. Default: \(1\) second.
  \item[\texttt{--max\_block\_draw\_trials}] Number of rejection-sampling
  attempts used to satisfy the block-start separation constraint. If no
  candidate passes after this many attempts, the next draw is accepted and
  the violation is recorded in the output. Default: \(30\).
  \item[\texttt{--num\_scrambles}] Number \(P\) of circular-shift null
  refits saved in this bag file. Default: \(4\).
  \item[\texttt{--min\_roll\_shift\_sec}] Two-sided exclusion zone around
  zero circular shift for each neuron. A shift is drawn from the interval
  excluding shifts smaller than this value and shifts closer than this value
  to a full wrap. Default: \(2\) seconds.
  \item[\texttt{--null\_A\_init}] Fresh \(A\) initialization mode for every
  null refit. Values: \texttt{scrambled\_data} (default) or \texttt{rand}.
  The real-time \(\hat A\) is deliberately not used to seed the null.
  \item[\texttt{--null\_B\_init}] Fresh \(\mathbf B\) initialization mode for
  every null refit. Values: \texttt{scrambled\_data} (default) or
  \texttt{rand}. The null \(\mathbf B\) is refit and saved as
  \texttt{B\_null}.
\end{description}

\paragraph{Inherited EM and M-step arguments.}
The Stage~(a) driver shares the ordinary PRISM-EM workhorse with
\texttt{prism\_EM\_train3c.py}.
\begin{description}
  \item[\texttt{--num\_em\_iters}] Number of outer EM iterations for the real
  time-ordered bag fit. Default: \(20\).
  \item[\texttt{--m\_epochs}] M-step Adam epochs per EM iteration. Default:
  \(2\). The null locked M-step uses the total count
  \(\texttt{num\_em\_iters}\times\texttt{m\_epochs}\); there is no separate
  \texttt{null\_m\_epochs}.
  \item[\texttt{--pgd\_iter}] PGD iterations per time bin in the E-step.
  Default: \(5\).
  \item[\texttt{--lr\_estep}] PGD learning rate in the E-step. Default:
  \(0.03\).
  \item[\texttt{--lambda2}] Temporal smoothness weight for the state
  probabilities. Default: \(2.0\).
  \item[\texttt{--lr\_mstep}] Adam learning rate for the M-step. Default:
  \(0.003\).
  \item[\texttt{--lr\_end\_factor}] Final learning-rate multiplier after
  linear decay. Default: \(0.1\).
  \item[\texttt{--lambda3}] L1 penalty weight on off-diagonal \(A\). Default:
  \(0.02\).
  \item[\texttt{--rho\_max}] Target spectral radius for soft correction of
  \(A\). Default: \(0.95\).
  \item[\texttt{--prescale\_m\_step\_4\_ArhoMax}] M-step batch interval for
  spectral-radius projection checks. Default: \(120\).
  \item[\texttt{--delay\_em\_iter\_4\_ArhoMax}] First real-fit EM iteration
  after which \(\rho_{\max}\) correction is active. Default: \(3\).
  \item[\texttt{--delay\_em\_iter\_4\_lrDecay}] First real-fit EM iteration
  after which learning-rate decay advances. Default: \(1\).
  \item[\texttt{--delay\_em\_iter\_4\_Aprune}] First real-fit EM iteration
  after which L1 proximal pruning is active. Default: \(1\).
  \item[\texttt{--batch\_size}] M-step minibatch size. Default: \(2048\).
  \item[\texttt{--init\_states}] Initial latent-state mode for the real fit:
  \texttt{data} (default) or \texttt{rand}.
  \item[\texttt{--init\_A}] Initial \(A\) mode for the real fit:
  \texttt{data} (default) or \texttt{rand}.
  \item[\texttt{--init\_B}] Initial \(\mathbf B\) mode for the real fit:
  \texttt{data} (default) or \texttt{rand}.
  \item[\texttt{--init\_A\_Tmax}] Maximum number of bins used for
  data-driven \(A\) initialization. Default: \(50000\).
  \item[\texttt{--decode\_dwell\_sec}] Dwell-time prior used in Viterbi
  decoding. Default: \(0.3\) seconds.
  \item[\texttt{--seed}] Base random seed. Default: \(42\); the effective
  seed is offset by \texttt{bag\_idx}.
  \item[\texttt{--fitName}] Accepted by the shared parser for ordinary EM
  compatibility, but Stage~(a) output naming is determined by
  \texttt{dataName} and \texttt{bag\_idx}.
  \item[\texttt{--verb}] Verbosity level. With the default verbosity, rank
  \(0\) reports runtime device information, real-fit elapsed time, and
  elapsed time for each scramble refit.
\end{description}

\paragraph{Saved Stage (a) contents.}
The file has the normal real-fit fields (\texttt{A\_hat}, \texttt{B\_hat},
\texttt{S\_hat}, \texttt{c\_hat}, training histories, etc.) plus
\texttt{A\_null}, \texttt{B\_null}, \texttt{roll\_shifts\_bin},
\texttt{block\_start\_bins}, \texttt{block\_start\_sec},
\texttt{boundary\_pair\_indices}, \texttt{block\_retry\_violations}, and
locked null M-step histories. The assembled bag spike train is not saved.
Stage~(a) metadata are grouped under \texttt{bagsFDR\_stageA}. This includes
the bag assembly controls, the circular-shift null-refit controls and
per-scramble initialization diagnostics, and a \texttt{real\_fit} sub-dict
holding the usual PRISM-EM \texttt{train}, \texttt{init\_A},
\texttt{init\_B}, \texttt{init\_state}, and \texttt{states\_recovery\_eval}
metadata for the real time-ordered fit on that bag.

\subsection{\texttt{prism\_EM\_FDR\_Bags\_agregate3c.py}: Stage (b)}

This program is a single CPU postprocessing step. It reads the Stage~(a)
files, performs the per-row null thresholding and cross-bag stability
selection, and writes one aggregate file in the same outer format expected
by \texttt{prism\_EM\_eval3c.py}. The filename contains
\texttt{agregate} with one ``g'' because that is the implemented script
name.

\paragraph{Input and output naming.}
The default input directory is
\[
\texttt{<basePath>/prismFDR}.
\]
For \texttt{num\_bags} \(=B\), the program expects the contiguous file set
\begin{verbatim}
<dataName>.bag000.prismFDRbag.npz
<dataName>.bag001.prismFDRbag.npz
...
<dataName>.bag{num_bags-1, zero padded}.prismFDRbag.npz
\end{verbatim}
Here \texttt{dataName} is the extended Stage~(a) name, i.e.\
\texttt{<originalDataName>\_<bagsTag>}.
The output stem is
\[
\texttt{<dataName>\_bags<num\_bags>},
\]
and the output file is written to
\[
\texttt{<basePath>/prismFit/<dataName>\_bags<num\_bags>.prismEM.npz}.
\]
For example, with \texttt{dataName=daleN100\_2ba29b\_c47b43} and
\texttt{num\_bags=5}, the output stem is
\[
\texttt{daleN100\_2ba29b\_c47b43\_bags5}.
\]
It can then be opened by
\begin{verbatim}
./prism_EM_eval3c.py \
  --basePath <basePath> \
  --dataName <dataName>_bags5 \
  -p a e f g
\end{verbatim}

\paragraph{Runtime arguments.}
\begin{description}
  \item[\texttt{--basePath}] Root run directory. The default input is
  \texttt{<basePath>/prismFDR}; the aggregate output is written to
  \texttt{<basePath>/prismFit}.
  \item[\texttt{--dataName}] Original dataset short name used by Stage~(a).
  This is the prefix of all bag files.
  \item[\texttt{--num\_bags}] Number of Stage~(a) bag files to read. The
  implementation expects bags numbered contiguously from \texttt{bag000}.
  \item[\texttt{--per\_bag\_quantile}] Quantile \(q\) used for the per-row
  null magnitude threshold \(\tau^{(b)}_i\). Default: \(0.99\).
  \item[\texttt{--sel\_prob}] Cross-bag selection-frequency threshold
  \(\pi_{\text{thr}}\). Default: \(0.7\).
  \item[\texttt{--fdr\_out\_dir}] Optional override for the Stage~(a) bag
  directory. Default: \texttt{<basePath>/prismFDR}.
  \item[\texttt{--verb}] Verbosity level. Default: \(1\).
\end{description}

\paragraph{Aggregate output schema.}
The output \texttt{A\_hat} is the final stable matrix: off-diagonal entries
are the selection-conditional bagged means for edges with
\(\pi_{ij}\ge\texttt{sel\_prob}\), and zero otherwise. With the default
\texttt{--sel\_prob 0.7}, this means retaining only edges present in at least
\(70\%\) of bags; the retained edge value is the unweighted mean over only
those bag instances where the edge was selected. Diagonal entries are
the all-bag mean diagonals. \texttt{A\_prune}, \texttt{neuron\_type}, and
\texttt{neuron\_Sedge} are recomputed from this aggregate \texttt{A\_hat}.
The output \texttt{B\_hat} is the mean across bags after aligning state rows
to bag~0 by the fitted \(\mathbf B\) vectors. The time-series fields
\texttt{S\_hat}, \texttt{c\_hat}, \texttt{S\_hat\_CL}, and the EM histories
are copied from reference bag~0 to preserve compatibility with the existing
plotting code; those plots should therefore be interpreted as reference-bag
diagnostics, while the connectivity fields are Stage~(b) aggregates.
Because \texttt{bag000} is chronological, these copied state assignments
are aligned to the original input timeline for simulation truth comparisons.
Stage~(b) metadata are grouped under \texttt{bagsFDR\_stageB}; Stage~(a)
metadata copied from reference bag~0 remain under \texttt{bagsFDR\_stageA}.

Additional Stage~(b) arrays include \texttt{A\_bag},
\texttt{B\_bag\_aligned}, \texttt{state\_permutation\_to\_ref},
\texttt{selected\_mask}, \texttt{selected\_in\_bag},
\texttt{selection\_count}, \texttt{selection\_frequency},
\texttt{A\_mean\_selected}, \texttt{A\_sd\_selected},
\texttt{A\_mean\_all}, \texttt{A\_sd\_all}, \texttt{row\_null\_tau\_bag},
\texttt{row\_null\_tau\_mean}, \texttt{row\_null\_mean},
\texttt{row\_null\_sd}, \texttt{z\_null}, and a compact selected-edge table
with \texttt{edge\_i}, \texttt{edge\_j}, \texttt{edge\_sel\_freq},
\texttt{edge\_A\_mean}, \texttt{edge\_A\_sd\_boot}, and
\texttt{edge\_z\_null}.

% ===============================================================
\section{End-to-End Summary}
% ===============================================================

\paragraph{Stage (a), run \(N_{\text{bag}}\) times in parallel (one node,
\(4\times\)A100 each).}
\begin{enumerate}
  \item Assemble bag \(b\): draw \(N_{\text{blk}}=10\) block starts in the
  time range \(R_{\text{time}}\), enforce
  \(\delta_{\text{blk}}^{\min}=1\,\text{s}\) separation by rejection sampling
  (up to \texttt{max\_block\_draw\_trials}=30 attempts per new block, then
  accept the next draw and record any violation), extract non-circular blocks
  of \(L_{\text{blk}}=60\,\text{s}\), concatenate, and record the block
  offsets and boundary pair indices.
  \item \(\textsf{full\_fit}\to(\hat A^{(b)},\hat S^{(b)},\hat{\mathbf B}^{(b)})\).
  \item For \(p=1,\dots,P\): draw independent per-neuron circular shifts in
  the two-sided zone
  \([\delta_{\text{roll}}^{\min,\text{bin}},T_b-\delta_{\text{roll}}^{\min,\text{bin}}]\),
  roll spikes, and
  \(\textsf{mstep\_locked}\to(A^{\text{null},(b,p)},\mathbf B^{\text{null},(b,p)})\)
  with the real pipeline's \(\lambda_3,\rho_{\max}\), decay active from epoch
  one, \(\hat S^{(b)}\) frozen and aligned as \(\hat S^{(b)}_{t-1}\) for pair
  \((Y_{t-1},Y_t)\), fresh \(A\) initialization controlled by
  \texttt{null\_A\_init}\(\in\{\texttt{scrambled\_data},\texttt{rand}\}\),
  and fresh \(\mathbf B\) initialization controlled by
  \texttt{null\_B\_init}\(\in\{\texttt{scrambled\_data},\texttt{rand}\}\).
  The number of locked M-step epochs and all other M-step hyperparameters
  are inherited from the real time-ordered fit configuration.
  \item Write one file for bag \(b\), named from \texttt{dataName} and
  \texttt{bag\_idx}, containing \(\hat A^{(b)}\),
  \(\{A^{\text{null},(b,p)}\}\), \(\{\mathbf B^{\text{null},(b,p)}\}\), \(\hat S^{(b)}\),
  realized block starts, boundary pair indices, block-start retry violations,
  realized shifts, fitted sparsity, and history summary. Do not save the
  assembled bag spike train.
\end{enumerate}

\paragraph{Stage (b), single postprocessing job.}
\begin{enumerate}
  \item Per bag and row, pool the \((N-1)P\) off-diagonal null weights, form
  the upper-\(q\) magnitude threshold \(\tau^{(b)}_i\), and select edges with
  \(|\hat A^{(b)}_{ij}|>\tau^{(b)}_i\); record \(q_\Lambda\).
  \item Compute selection frequencies \(\pi_{ij}\) and keep edges with
  \(\pi_{ij}\ge\pi_{\text{thr}}\) (or use the
  \texttt{target\_false\_edges}-derived \(\pi_{\text{thr}}\)).
  \item For surviving edges, compute the bagged mean \(\bar A_{ij}\),
  bootstrap SD \(\widehat{\mathrm{SD}}_{ij}\) (selection-conditional and
  all-bag), and effect size \(z_{ij}\) against the pooled row null.
  \item Emit the vetted edge list
  \((i,j,\pi_{ij},\bar A_{ij},\widehat{\mathrm{SD}}_{ij},z_{ij})\) plus
  diagnostics.
\end{enumerate}

\paragraph{Closing note.} The method's robustness comes from never trusting a
single fit: connectivity is asserted only when it survives data resampling
(bagging), exceeds a null built by the identical estimation machinery
(scramble calibration), and recurs across bags (stability selection). Its
two reported uncertainties --- the null spread and the across-bag bootstrap
SD --- are deliberately kept distinct, and their known limitations under
overlapping bags are stated rather than papered over, so that the resulting
edge list can be read with appropriate confidence on real data where no
ground truth is available to check it.

\end{document}
